\section{Introduction}

The parton distribution functions (PDFs) $f(x)$ \cite{Feynman:1973xc}
are related to matrix elements
of bilocal operators on the light cone $z^2=0$, which prevents
a straightforward calculation of these functions in the lattice
gauge theory
formulated in Euclidean space.
The usual way out is to calculate their moments.
However, recently, X. Ji \cite{Ji:2013dva} suggested a method
allowing to calculate PDFs as functions of $x$.
To this end, he proposes to use purely space-like
separations $z=(0,0,0,z_3)$.
Then one deals with quasi-PDFs $Q(y,p_3)$ describing sharing
of the $p_3$ hadron momentum component,
and
tending to PDFs $f(y)$
in the \mbox{$p_3 \to \infty$} limit.
The same method can be applied to distribution amplitudes (DAs).
The results of lattice calculations of
quasi-PDFs were reported
in Refs. \cite{Lin:2014zya,Chen:2016utp,Alexandrou:2015rja}
and of the pion quasi-DA in Ref. \cite{Zhang:2017bzy}.

In our recent papers \cite{Radyushkin:2016hsy,Radyushkin:2017gjd},
we have studied nonperturbative $p_3$-evolution of
quasi-PDFs and quasi-DAs
using the formalism of virtuality distribution functions
\cite{Radyushkin:2014vla,Radyushkin:2015gpa}.
We found that quasi-PDFs can be obtained from
the transverse momentum dependent
distributions (TMDs) ${\cal F} (x, k_\perp^2)$. We
built models for the nonperturbative evolution of
quasi-PDFs using simple models for TMDs.
Our results are in qualitative agreement with the
\mbox{$p_3$-evolution} patterns obtained in lattice calculations.

In the present paper, our first goal
is to develop a picture for
quasi-PDFs as hybrids of PDFs and
primordial
momentum distributions of partons in a hadron at rest.
As an intermediate step, we demonstrate that the connection between
TMDs and quasi-PDFs
\cite{Radyushkin:2016hsy} is a mere consequence
of Lorentz invariance. Then we show that,
when a hadron is moving, the parton $k_3$ momentum
comes from two sources. The motion
of the hadron as a whole gives
the $xp_3$ part, governed by the dependence of
the TMD ${\cal F} (x, \kappa^2)$ on its $x$ argument.
The remaining part $k_3-xp_3$
is governed by the dependence of the TMD on its second argument,
$\kappa^2$,
dictating the primordial
rest-frame
momentum distribution. The convolution nature of quasi-PDFs
results in a rather complicated pattern of their
\mbox{$p_3$-evolution,} necessitating rather large values $p_3 \sim 3$ GeV
for getting close to the PDF limit.

Thus, our second goal is to propose an alternative approach for
lattice PDF extraction. To this end, we introduce
{\it pseudo-PDFs} ${\cal P} (x, z_3^2) $ that generalize
the light-cone PDFs $f(x)$ onto spacelike intervals like
$z=(0,0,0,z_3)$. The pseudo-PDFs are Fourier transforms
of the \mbox{{\it Ioffe-time} \cite{Ioffe:1969kf}} {\it distributions} \cite{Braun:1994jq}
${\cal M} (\nu, z_3^2)$
that are basically given by generic matrix elements like $
\langle p | \phi(0) \phi (z)|p \rangle $
written as functions of $\nu = p_3 z_3$ and $z_3^2$.
Unlike quasi-PDFs, the pseudo-PDFs have the ``canonical'' $-1 \leq x \leq 1$ support
for all $z_3^2$. They tend to PDFs when $z_3\to 0$,
showing in this limit a usual perturbative evolution
with $1/z_3$ serving as an evolution parameter.
Finally, we discuss how these properties of pseudo-PDFs may be used
for extraction of PDFs on the lattice.
